%=============================================================================
% Section 3: Why Minimal Execution Matters
%=============================================================================
\section{Why Minimal Execution Matters}

Minimality is treated as a system constraint rather than a stylistic preference. Embedding policy logic into the execution runtime introduces state coupling, reduces portability, and increases cross-layer failure propagation. Minimality is therefore defined as follows:

\begin{quote}
\textbf{Minimality Constraint:} The execution layer MUST implement deterministic capability invocation only, and MUST NOT perform policy-driven decisions such as provider selection, strategic retry, or workflow-level fallback.
\end{quote}

Under this constraint, runtime semantics are specified through core responsibilities and explicit non-responsibilities.

\subsection{Core Responsibilities}

The runtime serves as a standardized capability execution engine. Its responsibilities are intentionally limited to six core functions:

\textbf{1. Capability Invocation:} Execute a capability defined in the AI Protocol. The runtime receives standardized capability requests and invokes the appropriate execution path.

\textbf{2. Protocol Translation:} Convert protocol-level requests into provider-specific API calls. This translation layer enables applications to work with standardized capability definitions while the runtime handles provider-specific implementations.

\textbf{3. Result Normalization:} Return standardized responses to the caller. Regardless of which provider executed the capability, the runtime normalizes results into a consistent format.

\textbf{4. Connection Handling:} Manage network communication with providers. This includes connection pooling, timeout management, and network-level error handling.

\textbf{5. Error Propagation:} Report execution errors without implementing application-level retry strategies. The runtime distinguishes between execution errors (which it reports) and retry logic (which belongs in the orchestration layer).

\textbf{5.1. Bounded Retry Semantics (Optional):} For transient network failures (e.g., connection timeouts, 429 rate limits), the runtime MAY implement limited micro-retries---typically 1--2 immediate retries with exponential backoff at the network layer. This differs from application-level retry strategies in that:
\begin{itemize}
    \item Micro-retries are bounded and automatic
    \item They address only transient network-level failures
    \item They do not implement policy-driven retry logic (e.g., provider fallback)
    \item Total retry duration must remain within the execution timeout budget
\end{itemize}

This choice reflects production constraints: transient network failures are common, and surfacing every transient event to the orchestration layer introduces avoidable latency and cross-layer communication overhead.

\textbf{6. Session Transport:} Forward session context without interpreting application policy. The runtime transports contextual information required for capability execution but does not interpret or modify policy decisions.

\subsection{What the Runtime Does NOT Do}

To maintain its minimal design, the runtime explicitly avoids:
\begin{itemize}
    \item \textbf{Agent Planning:} Deciding what sequence of capabilities to invoke
    \item \textbf{Workflow Orchestration:} Managing complex multi-capability workflows
    \item \textbf{Application Logic:} Implementing business logic specific to applications
    \item \textbf{Provider Selection:} Choosing which provider should execute a capability
    \item \textbf{Routing Policies:} Implementing cost optimization or latency-based routing
    \item \textbf{Application-Level Retry Strategies:} Determining when and how to retry failed executions across providers or workflows
\end{itemize}

These responsibilities belong in the orchestration and agent layers, maintaining clear separation of concerns.

\subsection{Stateless Execution Design}

The runtime operates as a stateless execution engine. Each capability invocation is independent, with any necessary state maintained by the calling agent or orchestration layer. This design enables:
\begin{itemize}
    \item \textbf{Horizontal Scalability:} Multiple runtime instances can be deployed without state synchronization
    \item \textbf{Simplified Debugging:} Each execution is self-contained and traceable
    \item \textbf{Fault Isolation:} Failures in one execution do not affect others
\end{itemize}

\subsection{Asynchronous and Long-Running Tasks}

Modern AI capabilities extend beyond synchronous text generation to include long-running tasks such as video generation, large-scale search indexing, and multi-step document processing. The execution layer accommodates these patterns while maintaining its minimal design:

\textbf{Streaming Responses:} For capabilities supporting Server-Sent Events (SSE) or WebSocket connections, the runtime:
\begin{itemize}
    \item Establishes the streaming connection on behalf of the caller
    \item Forwards stream chunks without buffering the entire response
    \item Reports stream errors immediately without buffering delays
    \item Allows the orchestration layer to handle stream cancellation
\end{itemize}

\textbf{Polling-Based APIs:} For asynchronous capabilities that return job IDs, the runtime:
\begin{itemize}
    \item Executes the initial API call and returns the job identifier
    \item Delegates polling logic to the orchestration layer (which has policy context)
    \item Provides a separate \texttt{check\_status} operation for poll queries
    \item Does not implement automatic polling or status change notifications
\end{itemize}

\textbf{Design Rationale:} By keeping long-running task state management in the orchestration layer, the runtime remains stateless and does not require background job tracking mechanisms. This preserves horizontal scalability and simplifies fault recovery.

\subsection{Security Considerations}

As the action interface between AI cognition and external capabilities, the execution layer occupies a critical position in the security model. While the runtime intentionally avoids business logic, it should implement baseline security measures:

\textbf{Input Validation:} The runtime validates capability requests against the protocol schema before execution. All \texttt{input\_schema} definitions MUST conform to JSON Schema 2020-12 or later, ensuring interoperability and standardized validation across implementations. The runtime performs schema validation at invocation time, catching structural errors before they reach providers.

\textbf{PII Handling (Optional):} Organizations may configure PII stripping or redaction at the runtime layer. This serves as a defense-in-depth measure against prompt injection attacks that attempt to exfiltrate sensitive data. However, this is a configurable option---primary PII protection remains the responsibility of the agent and application layers.

\textbf{Rate Limiting:} The runtime may implement local rate limiting to prevent runaway capability invocations, protecting both the runtime infrastructure and downstream providers.

\textbf{Audit Logging:} All capability invocations should be logged for security auditing, including the capability name, provider, timestamp, and execution status (success/failure).

To enable observability for the orchestration layer, the runtime SHOULD include the following metrics in standardized response metadata:
\begin{itemize}
    \item \texttt{retry\_count}: Number of micro-retries performed during execution (if applicable)
    \item \texttt{translation\_latency\_ms}: Time spent in protocol translation
    \item \texttt{execution\_latency\_ms}: Time spent in actual provider API call
\end{itemize}

These metrics allow the Contact Layer to assess provider stability and make informed routing decisions.

\textbf{Security Boundary:} The execution layer is not a security perimeter. It operates within a trusted environment where the orchestration layer has already authenticated the agent and authorized the capability request. The runtime's security role is defense-in-depth, not primary access control.
